\section{定積分では端点の項を先に見る}\label{節:定積分の部分積分}

不定積分を求めてから上下端を代入してもよいが，最初から定積分の公式を使えば，端点で消える項がすぐに分かる。公式
\[
 \int_a^b uv'\dd x=[uv]_a^b-\int_a^b u'v\dd x
\]
では，残る積分だけでなく，端点の $u(a),u(b),v(a),v(b)$ も計算の見通しに関わる。積のどちらかが $0$ ならば，その端点の寄与はなくなる。

\begin{example}{端点で消える積を確かめる}
$\int_0^{\frac\pi2}x\sin x\dd x$ を求めよ。
\end{example}

$u=x$，$v'=\sin x$ と選ぶと，$u'=1$，$v=-\cos x$ なので，
\[
\begin{aligned}
 \int_0^{\frac\pi2}x\sin x\dd x
 &=[-x\cos x]_0^{\frac\pi2}
   +\int_0^{\frac\pi2}\cos x\dd x\\
 &=0+[\sin x]_0^{\frac\pi2}=1.
\end{aligned}
\]
上端では $\cos\frac\pi2=0$，下端では $x=0$ のため，端点の項はともに $0$ になった。二つの端点で消える理由は異なるので，一方だけを見て項全体を消しているわけではない。

また，積分される関数はこの区間で非負だから，求めた値が正であることは自然である。仮に $-1$ になったら，$\sin x$ の原始関数の負号か，部分積分の前の負号を確かめるべきだと分かる。

\begin{example}{部分積分の後に割り算を使う}
$\int_0^1x\log(1+x)\dd x$ を求めよ。
\end{example}

$\log(1+x)$ を微分すると $\frac1{1+x}$ になる。$x$ の原始関数は $\frac{x^2}{2}$ だから，残る積分は有理関数になる。そこで
\[
 u=\log(1+x),\qquad v'=x
\]
と選び，
\[
 \int_0^1x\log(1+x)\dd x
 =\left[\frac{x^2}{2}\log(1+x)\right]_0^1
  -\frac12\int_0^1\frac{x^2}{1+x}\dd x
\]
を得る。第一項は $\frac12\log2$ である。第二項では，分子の次数が分母より高いので，\chapterref{章:基本積分}で学んだ割り算を使う。
\[
 x^2=(1+x)(x-1)+1,
 \qquad
 \frac{x^2}{1+x}=x-1+\frac1{1+x}.
\]
したがって
\[
\begin{aligned}
 \int_0^1\frac{x^2}{1+x}\dd x
 &=\left[\frac{x^2}{2}-x+\log(1+x)\right]_0^1\\
 &=-\frac12+\log2.
\end{aligned}
\]
これを戻すと
\[
 \int_0^1x\log(1+x)\dd x
 =\frac12\log2-\frac12\left(-\frac12+\log2\right)
 =\frac14.
\]
部分積分の直後には対数が端点に残り，別の積分にも対数が現れたが，最後に差を取ると打ち消し合った。途中の式の見た目だけで，答えにも対数が必ず残ると決めることはできない。

ここまでで，式の整理，置換積分，部分積分という三つの道具を学んだ。整式の割り算や三角恒等式で和へ直せるなら，まずその変形を試せる。合成関数と内側の微分がそろっていれば置換を考え，微分すると簡単になる因子と，積分できる因子の積なら部分積分を考える。ある操作で何が簡単になり，何が残るかを具体的に見ることが，選択の根拠になる。

ただし，定積分の値だけが欲しいときは，原始関数を最後まで求める必要がない場合もある。区間の左右を入れ換えたときに式がどう変わるか，正負の部分がどう対応するかを使えば，値が決まることがある。次の\chapterref{章:定積分の対称性}では，被積分関数と積分区間を一緒に見て，その対応を計算へ生かす。

